\section{Technical Overview}
\label{sec:overview}

Our main idea is to reduce a length-\(N\) vector relation over a small
finite field \(\Field\) to a length-\(m\) vector relation over
\(\ExtensionField\), where \(N=m\kappa\) and
\([\ExtensionField:\Field]=\kappa\).
We achieve this through \emph{polynomial packing}: encode the original
relation as a low-degree identity on length-\(m\) vector polynomials,
then evaluate this identity at a random extension-field point.
The prover performs the row-polynomial multiplications over \(\Field\)
and invokes an extension-field proof only on the resulting shorter
relation. A packed polynomial commitment authenticates its operands
against the original witness.

\subsection{Polynomial Packing of Vector Relations}
\label{sec:dualPhase}

Assume \(m,\kappa\) are powers of two and \(\kappa\geq2\).
Arrange the committed vectors \(\bm a,\bm b\in\Field^N\) as
\(\kappa\) columns of length \(m\):
\[
    \bm A=\concat{\bm a^{(1)},\ldots,\bm a^{(\kappa)}},
    \qquad
    \bm B=\concat{\bm b^{(1)},\ldots,\bm b^{(\kappa)}}
    \in\Field^{m\times\kappa}.
\]
Each row contains \(\kappa\) base-field elements. We encode that row
by one polynomial of degree below \(\kappa\), so the whole matrix
becomes a length-\(m\) vector of polynomials.
For a basis
\(\bm\phi(X)=(\phi_1(X),\ldots,\phi_\kappa(X))^\top\), with
\(\phi_j\in\Field[X]\) of degree below \(\kappa\), define:
\begin{equation}
    \bm g_A(X)=\bm A\bm\phi(X),
    \qquad
    \bm g_B(X)=\bm B\bm\phi(X)
    \in\Field[X]^m.
\label{eq:overview-packing}
\end{equation}
Evaluating a polynomial vector replaces each row of \(\kappa\)
base-field elements by one extension-field element.
At \(\ext x\in\ExtensionField\), for example:
\[
    \bm A\in\Field^{m\times\kappa}
    \quad\longmapsto\quad
    \bm g_A(X)
    \quad\xrightarrow{X=\ext x}\quad
    \ext{\bm a}=\bm g_A(\ext x)\in\ExtensionField^m.
\]
To turn this compression into a reduction of the \emph{relation},
we proceed in two phases.

\paragraph{Encode the relation as a polynomial identity.}
We choose the basis so that the original claim is expressed by
coefficients or interpolation values of a polynomial expression in
\(\bm g_A\) and \(\bm g_B\).
The prover supplies \(O(\kappa)\) auxiliary coefficients for the
remaining terms; their prescribed form prevents them from changing
the encoded claim.
The inner- and Hadamard-product constructions below instantiate this
step with monomial and Lagrange bases, respectively.

\paragraph{Evaluate to obtain the shorter relation.}
After the auxiliary coefficients are fixed, the verifier samples an
extension-field challenge. The prover evaluates the vector polynomials,
while the verifier evaluates the short auxiliary polynomials.
The identity becomes a relation on length-\(m\) extension-field
vectors, with a verifier-computable target.
For our constructions, an incorrect encoded claim leaves a
nonzero residual of degree \(O(\kappa)\), after clearing negative
powers, giving reduction error \(O(\kappa/|\ExtensionField|)\).
The resulting relation is a single inner product: its proof costs
\(O(m)\) extension-field operations and leaves two multilinear-evaluation
(MLE) claims on the packed operands~\cite{asiacrypt/CCYY25}.
The commitment scheme below reduces these claims to openings of the
original committed witness.

\subsection{Instantiating the Packed Reduction}

\paragraph{Inner products.}
For \(c=\langle\bm a,\bm b\rangle\), we want all matching-column
products to contribute to the same coefficient. Opposite monomial
powers achieve exactly this: choose
\(\bm\phi(X)=(1,X,\ldots,X^{\kappa-1})^\top\) and pair
\(\bm g_A(X)\) with \(\bm g_B(X^{-1})\).
For two blocks:
\[
\begin{aligned}
 &\langle\bm a^{(1)}+X\bm a^{(2)},
          \bm b^{(1)}+X^{-1}\bm b^{(2)}\rangle\\
 &\quad=
   \langle\bm a^{(1)},\bm b^{(1)}\rangle
   +\langle\bm a^{(2)},\bm b^{(2)}\rangle\\
 &\qquad\quad
   +X\langle\bm a^{(2)},\bm b^{(1)}\rangle
   +X^{-1}\langle\bm a^{(1)},\bm b^{(2)}\rangle.
\end{aligned}
\]
The constant coefficient is precisely the claimed inner product;
every cross-column term has a nonzero exponent. Thus, for arbitrary
\(\kappa\), the claim holds if and only if there exist
\(h_+,h_-\in\Field[X]\), each of degree at most \(\kappa-2\), satisfying:
\begin{equation*}
    \langle\bm g_A(X),\bm g_B(X^{-1})\rangle
    =c+Xh_+(X)+X^{-1}h_-(X^{-1}).
\end{equation*}
No choice of \(h_+\) or \(h_-\) can compensate for an incorrect \(c\).
Once the prover sends these polynomials, the verifier samples
\(\ext x\in\ExtensionField^\times\).
Writing \(\ext{\bm a}=\bm g_A(\ext x)\) and
\(\ext{\bm b}=\bm g_B(\ext x^{-1})\), it remains to prove:
\[
    \langle\ext{\bm a},\ext{\bm b}\rangle
    =c+\ext xh_+(\ext x)+\ext x^{-1}h_-(\ext x^{-1}).
\]
This is a length-\(m\) inner-product claim with a verifier-computable
target. Its auxiliary polynomials contain only sums of base-field
cross-column products.

\paragraph{Hadamard products.}
The inner-product encoding retains the sum of matching-column
products. A Hadamard relation also distinguishes their locations.
For a committed \(\bm C\in\Field^{m\times\kappa}\), we first project
\(\bm C=\bm A\circ\bm B\) using random tensor weights
\(\ext{\bm l}\in\ExtensionField^m\) and
\(\ext{\bm r}\in\ExtensionField^\kappa\).
Each is a Kronecker product of two-entry MLE vectors, so this is a
random evaluation of the residual's multilinear extension.
A false matrix equality survives with probability at most
\(\log_2N/|\ExtensionField|\).
We authenticate \(\ext c=\ext{\bm l}^{\top}\bm C\ext{\bm r}\) and prove:
\[
    \ext c=\ext{\bm l}^{\top}(\bm A\circ\bm B)\ext{\bm r}.
\]
The prover supplies claimed column products
\(\ext c_j=\langle\ext{\bm l},
\bm a^{(j)}\circ\bm b^{(j)}\rangle\), and the verifier checks
\(\ext c=\ext{\bm c}^{\top}\ext{\bm r}\).

To verify these column claims together, we encode each at its own
interpolation node. Choose \(\kappa\) distinct nodes
\(\mathcal H=\{\alpha_1,\ldots,\alpha_\kappa\}\subseteq\Field\)
and the Lagrange basis satisfying
\(\psi_j(\alpha_i)=1\) for \(i=j\) and \(0\) otherwise.
Use \(\bm\phi=(\psi_1,\ldots,\psi_\kappa)^\top\) in
Eq.~\eqref{eq:overview-packing}, and set
\(g_C(X)=\sum_{j=1}^{\kappa}\ext c_j\psi_j(X)\) and
\(\Psi(X)=\prod_{j=1}^{\kappa}(X-\alpha_j)\).
At node \(\alpha_j\), the packed vectors recover column \(j\).
Consequently, all column claims hold exactly when the residual
vanishes on \(\mathcal H\), or equivalently:
\begin{equation*}
    \ext{\bm l}^{\top}
      \bigl(\bm g_A(X)\circ\bm g_B(X)\bigr)-g_C(X)
    =\Psi(X)h(X),
    \qquad \deg h\leq\kappa-2,
\end{equation*}
where \(h\in\ExtensionField[X]\).
Here the factor \(\Psi\) prevents the auxiliary polynomial from
changing any claimed node value, just as the nonzero powers protected
the constant coefficient in the inner-product reduction.
After \(\ext{\bm c}\) and \(h\) are fixed, the verifier samples
\(\ext x\in\ExtensionField\setminus\mathcal H\).
With \(\ext{\bm a}=\bm g_A(\ext x)\) and
\(\ext{\bm b}=\bm g_B(\ext x)\), the reduced claim is:
\[
    \langle\ext{\bm a},\ext{\bm b}\circ\ext{\bm l}\rangle
    =g_C(\ext x)+\Psi(\ext x)h(\ext x).
\]

Thus both instances reduce the original length-\(N\) relation to a
length-\(m\) inner product over \(\ExtensionField\). The basis and
auxiliary polynomials differ; the subsequent extension-field proof
and commitment interface are the same.

\subsection{Prover Complexity}
\label{sec:highLevelEfficiency}

Polynomial packing confines the extension-field inner-product proof
to \(m=N/\kappa\) entries. To obtain a net saving, the preceding
polynomial computation must also exploit the small witness field.
We therefore multiply row polynomials before evaluating them over
\(\ExtensionField\).
For inner products, this computes \(h_+\) and \(h_-\) entirely over
\(\Field\). For Hadamard products, we first form the base-field row
products, then accumulate their coefficients with weights
\(\ext{\bm l}\) and divide the short residual by \(\Psi\).
Thus extension-field weights do not turn the row multiplications
into full extension-field multiplications.

With precomputed bases, these steps cost
\(O(m\kappa^2)=O(\kappa N)\) base-field operations and \(O(N)\)
mixed-field operations. Evaluating the packed vectors also costs
\(O(N)\) mixed-field operations, followed by \(O(m)\) extension-field
operations for the inner-product reduction.
Let \(\bb,\be,\ee\) denote the respective multiplication costs.
Under schoolbook arithmetic,
\(\be=\kappa\bb\) and \(\ee=\kappa^2\bb\).
For the bases above and \(m\geq\kappa\), the reduction costs:
\[
    O(\kappa N)\,\bb+O(N)\,\be+O(N/\kappa)\,\ee
    =O(\kappa N)\,\bb.
\]
Row products can be accumulated directly into the auxiliary
coefficients. Storing the witness, a constant number of length-\(m\)
extension-field vectors, and \(O(\kappa)\) auxiliary coefficients
therefore takes \(O(N+\kappa^2)=O(N)\) base-field coordinates.
In comparison, the uncompressed extension-field proof costs
\(O(N\kappa^2)\) base-field operations and \(O(N\kappa)\) coordinates.
These bounds exclude the commitment scheme.

\subsection{Authenticating Packed Evaluations}
\label{sec:highlevelPCS}

To complete the length reduction, we must authenticate evaluations of
the packed operands \(\bm A\ext{\bm\rho}\) using a commitment to
the original witness. Since the packing weights depend on the verifier's
challenge, we use a fixed packing for the commitment and derive the
required evaluations from it.
Fix an \(\Field\)-basis
\(\ext{\bm\beta}=(\ext\beta_1,\ldots,\ext\beta_\kappa)^\top\)
of \(\ExtensionField\), and commit to
\(\ext{\bm a}=\bm A\ext{\bm\beta}\in\ExtensionField^m\).
Since the entries of \(\bm A\) lie in \(\Field\), this packing is
injective and binds every column of \(\bm A\).

\paragraph{From extension-field claims to base-field coordinates.}
A terminal claim has the form
\(\ext u=\langle\bm A\ext{\bm\rho},\ext{\bm\chi}\rangle\),
where \(\ext{\bm\rho}\in\ExtensionField^\kappa\) gives the public
packing weights and \(\ext{\bm\chi}\) is a public tensor query.
For the weighted Hadamard operand, absorb \(\ext{\bm l}\) into the query;
coordinatewise multiplication preserves its tensor structure.
The prover supplies \(\ext{\bm v}=\bm A^\top\ext{\bm\chi}\), and
the verifier checks \(\ext u=\ext{\bm\rho}^{\top}\ext{\bm v}\).

Binding \(\bm A\) does not yet authenticate these column evaluations.
Unlike the witness entries, the entries of \(\ext{\bm v}\) lie in
\(\ExtensionField\), so combining them with \(\ext{\bm\beta}\) is not
injective. We restore injectivity by expressing the query and claimed
values in base-field coordinates:
\[
    \ext{\bm\chi}=\bm\Xi\ext{\bm\beta},
    \qquad
    \ext{\bm v}=\bm V\ext{\bm\beta},
    \qquad
    \bm\Xi\in\Field^{m\times\kappa},
    \quad \bm V\in\Field^{\kappa\times\kappa}.
\]
The claim becomes a base-field matrix identity. Packing its columns
now preserves the entire identity and exposes the committed vector:
\begin{equation*}
    \bm V=\bm A^\top\bm\Xi
    \quad\Longleftrightarrow\quad
    \ext{\bm\beta}^{\top}\bm V
    =\ext{\bm a}^{\top}\bm\Xi.
\end{equation*}
After \(\bm V\) is fixed, a uniform
\(\ext{\bm r}\in\ExtensionField^\kappa\) reduces this row identity
to one scalar claim, with error at most \(1/|\ExtensionField|\):
\[
    \ext{\bm\beta}^{\top}\bm V\ext{\bm r}
    =\langle\ext{\bm a},\bm\Xi\ext{\bm r}\rangle.
\]
An extension-field inner-product proof leaves an MLE opening of
\(\ext{\bm a}\) and a query-side claim
\(\ext w=\ext{\bm\xi}^{\top}\bm\Xi\ext{\bm r}\), where
\(\ext{\bm\xi}\in\ExtensionField^m\) is its terminal tensor weight.

\paragraph{Base-field challenges make the query check decodable.}
Although \(\ext{\bm\chi}\) is a tensor query, its coordinate columns
in \(\bm\Xi\) need not be tensor vectors.
Moreover, contracting with the extension-field weight
\(\ext{\bm\xi}\) produces extension-field values that cannot be
recovered as the coordinates of a single packed evaluation.
A base-field terminal weight would avoid this problem: contraction
then commutes with taking base-field coordinates.
Decompose \(\ext{\bm\xi}=\bm U\ext{\bm\beta}\), with
\(\bm U\in\Field^{m\times\kappa}\).
The prover sends a claimed product
\(\bm W=\bm U^\top\bm\Xi\in\Field^{\kappa\times\kappa}\);
the verifier checks \(\ext w=\ext{\bm\beta}^{\top}\bm W\ext{\bm r}\).
Since the matrix product is over \(\Field\), we can check it using
base-field challenges.
After \(\bm W\) is fixed, the verifier independently samples uniform
\(\bm s_{(i)},\bm t_{(i)}\in\Field^\kappa\) in each of \(k\)
repetitions; the prover proves:
\[
    \bm s_{(i)}^\top\bm W\bm t_{(i)}
    =\langle\bm U\bm s_{(i)},\bm\Xi\bm t_{(i)}\rangle.
\]
The inner-product reduction also uses base-field challenges, so its
terminal tensor weight \(\bm l_{(i)}\) lies in \(\Field^m\).
The two rows needed to check its terminal values are now recoverable
as the base-field coordinates of:
\[
    \langle\bm l_{(i)},\ext{\bm\xi}\rangle
       =(\bm l_{(i)}^\top\bm U)\ext{\bm\beta},
    \qquad
    \langle\bm l_{(i)},\ext{\bm\chi}\rangle
       =(\bm l_{(i)}^\top\bm\Xi)\ext{\bm\beta}.
\]
Both left-hand sides are inner products of tensor vectors, each
computable in \(O(\log m)\) extension-field operations.
Combining the recovered rows with \(\bm s_{(i)}\) and
\(\bm t_{(i)}\) checks the two terminal values without constructing
\(\bm U\) or \(\bm\Xi\).

Independent repetition amplifies soundness over the smaller challenge
field.
One repetition has error at most \((2\log_2m+2)/|\Field|\),
including the matrix projection and inner-product reduction.
For security parameter \(\lambda\), choose \(k\) so that
\(((2\log_2m+2)/|\Field|)^k\leq2^{-\lambda}\), assuming
\(|\Field|>2\log_2m+2\).
The prover computes \(\bm W\) in \(O(\kappa N)\) base-field operations
and spends \(O(kN)\) on the repetitions.
Processing them sequentially uses \(O(N+\kappa^2)\) base-field
coordinates, with no additional witness commitment.
